\documentclass[10pt,a4paper]{amsart}
\usepackage[margin=19mm]{geometry}
\usepackage{amsmath,amssymb,mathtools}
\usepackage[scr=rsfs,cal=euler]{mathalfa}
\usepackage{xfrac}
\usepackage[unicode,hidelinks,bookmarks=false]{hyperref}
\newcommand{\ie}{i.e.\ }
\newcommand{\cf}{cf.\ }
\newcommand{\resp}{respectively\ }
\newcommand{\Subsection}{\subsection}
\providecommand{\nobreakdash}{\nobreak\hbox{-}}
\setlength{\emergencystretch}{2em}
\multlinegap0pt
\usepackage{amssymb}
\usepackage{xcolor}
\usepackage{float}
\usepackage{multirow}
\newtheorem{lemma}{Lemma}
\title{\textbf{FEKAN: Feature-Enriched Kolmogorov-Arnold Networks}}
\author{Sidharth S. Menon, Ameya D. Jagtap}
\date{Source: arXiv:2602.16530v1 (CC BY 4.0)}
\begin{document}
\maketitle

\noindent\emph{Source and licence.} \textbf{FEKAN: Feature-Enriched Kolmogorov-Arnold Networks}. Version arXiv:2602.16530v1 (CC BY 4.0), distributed by arXiv under the Creative Commons Attribution 4.0 International licence, http://creativecommons.org/licenses/by/4.0/, as recorded in the arXiv licence metadata for this exact version and shown on the abstract page https://arxiv.org/abs/2602.16530v1. Full licence notice: \texttt{licenses/arxiv-2602.16530-CC-BY-4.0.txt}.\par\smallskip

\noindent\textit{Editorial scope.} Counted unit: the article's lemma lem:complexity-reduction (source0.tex lines 1815--1824). The article's complete original proof of it is delivered with it (source0.tex lines 1826--1897, one \texttt{proof} environment of 72 lines). No mathematics is written, repaired or re-derived: the statement and the proof are the article's own lines, in the article's own notation, and no retained line is substituted or re-flowed.  No further source-local context is needed: every cross-reference of every retained line resolves inside the two delivered spans.  Nothing was omitted from any retained span, so the package carries no omission record.  No retained line contains a citation, so the wrapper carries no bibliography and no bibliography entry.  No retained line contains a figure environment, an image-inclusion command or drawing code, so no image is delivered and the package declares no asset dependency.  Editorial formatting only, in addition: an AMS \texttt{amsart} preamble that restates the article's own declarations and macros used by the retained lines, namely the environments \texttt{lemma} on separate counters, together with the standard packages those lines need.  The retained environments print in this document as Lemma 1 in order of appearance, whereas the article prints the same items with the numbering of its own full document.  The complete native source of the article as distributed in the arXiv e-print for this exact version is bundled unchanged as \texttt{sources/194893/source0.tex}.
\begin{lemma}
\label{lem:complexity-reduction}
There exist continuous features \(u_j\) and a continuous target 
\(F:[0,1]^n \to \mathbb{R}\) such that, for some \(\varepsilon>0\),
\[
\mathrm{Comp}_{n,m}(F,\varepsilon)
\ <\
\mathrm{Comp}_{n}(F,\varepsilon).
\]
\end{lemma}

\begin{proof}
We construct a function that is extremely difficult to approximate without the
added nonlinear feature, but becomes trivially representable once that feature 
is included.

\paragraph{A highly oscillatory hidden feature:}
Define
\[
u_1(\mathbf{x})=\sin\!\big(N(x_1+\cdots+x_n)\big), \qquad N\gg 1,
\]
and let the target function be
\[
F(\mathbf{x})=g(u_1(\mathbf{x})),
\]
where $g$ is any continuous function.  
The key point is that $u_1$ oscillates $O(N)$ times along the direction 
$(1,1,\dots,1)$; hence so does $F$.

\paragraph{ Representation when the nonlinear feature is available:}
In the feature-augmented model, we allow an additional feature $\psi_{n+1}$ 
depending on all coordinates.  Define
\[
\psi_{n+1}(\mathbf{x}) = u_1(\mathbf{x}), 
\qquad
\psi_1 = \cdots = \psi_n = 0,
\]
and choose the outer functions
\[
\phi_1 = g,
\qquad 
\phi_2=\cdots=\phi_K=0.
\]
Then the composed representation satisfies
\[
\phi_1\big(\psi_{n+1}(\mathbf{x})\big) = g(u_1(\mathbf{x})) = F(\mathbf{x}).
\]
Thus the function is represented \emph{exactly} with a single outer term:
\[
\mathrm{Comp}_{n,m}(F,0) = 1.
\]

\paragraph{Difficulty when the nonlinear feature is unavailable:}
Now restrict to the diagonal
\[
(x_1,\dots,x_n)=(t,\dots,t).
\]
Along this line,
\[
F(t,\dots,t)=\sin(N n t),
\]
which oscillates $O(N)$ times as $t$ varies.

However, any function in the class $\mathcal{R}_n$ (the model without added
features) is a sum of $K$ univariate components applied to the individual
coordinates.  Along the diagonal, each such component reduces to a function of
\emph{one variable, but with no access to the high-frequency combination}
$t \mapsto Nt$ except through $t$ itself.  
A sum of only $K\ll N$ univariate functions cannot approximate a function that
oscillates $O(N)$ times; hence any such approximant incurs uniform error at 
least some fixed $\varepsilon>0$.  Therefore,
\[
\mathrm{Comp}_n(F,\varepsilon) \;\ge\; C > 1,
\]
where $C$ depends on $N$ but is uniformly bounded away from $1$ for 
$K\ll N$.

\paragraph{Conclusion:}
The function $F$ requires many components to approximate without the nonlinear
feature, but becomes representable \emph{exactly} with a single term once the
feature $u_1$ is added.  Hence representation complexity decreases strictly when 
nonlinear features are included.
\end{proof}

\end{document}
